% !PRINTED 158
\weberpage{158}

\begin{center}
\S.\,38.

\gesperrt{Substitution of the Ratios.}
\end{center}

We have already seen above that two substitutions of the same dimension which are similar to a third are similar to one another. Accordingly we can unite all mutually similar substitutions of dimension $n$ into one class, and every substitution can be placed in one and only one such class. The individual substitutions of a class are called the \gesperrt{representatives} of the class, and every class is completely determined by any one of its representatives.

If $A$ is similar to $A'$ and $B$ similar to $B'$, then $AB$ is also similar to $A'B'$.

If, then, we denote by $\mathfrak{A}, \mathfrak{B}$ the classes to which $A, A'$ and $B, B'$ belong, one always arrives at the same class, whatever representatives from $\mathfrak{A}$ and from $\mathfrak{B}$ one may compose. This class, which is represented by $AB$ or $A'B'$, we therefore call composed of $\mathfrak{A}$ and $\mathfrak{B}$, and we denote it by $\mathfrak{A}\mathfrak{B}$. In this composition the same rules hold as in the composition of the substitutions themselves, and the totality of the classes therefore likewise forms a group.

This conception of the substitutions and their composition is always expedient when, as for example in projective geometry, only the ratios of the variables matter, and we therefore designate the substitution classes as \gesperrt{substitutions of the ratios.}

We shall, however, denote these substitutions of the ratios in the same way as the other substitutions, namely each class by a representative, whereby a misunderstanding will not easily arise. The representative can be chosen from very different points of view.

It is often advisable to take it so that the determinant is $= 1$. But in that case, for a substitution of dimension $n$, the multiplicative substitution $(\mu)$ is determined only up to an $n$th root of unity. If the representatives are so chosen, this property is preserved under composition, and these representatives of the classes therefore also form among themselves a substitution group.
